\chapter{Guía práctica de la aplicación}
\label{chap:guia}

\section{Un primer recorrido reproducible}
QSB ofrece una aplicación macOS y una herramienta de terminal llamada
\texttt{qsb}. Ambas comparten los modelos, validadores y solvers de
\texttt{QSBCore}. La interfaz actual usa etiquetas en inglés; este manual
conserva esas etiquetas para que puedan localizarse en pantalla.

Para comenzar, abra \path{manual/examples/lp_sample.json} mediante
\textbf{Model > Open Model...}. El modelo maximiza $50X_1+60X_2$ con
$2X_1+3X_2\leq180$, $3X_1+2X_2\leq150$ y variables no negativas.
\begin{enumerate}
 \item En \textbf{Definition}, compruebe nombres, coeficientes, relaciones,
 límites y tipos de variables. Aquí las variables son continuas.
 \item Pulse \textbf{Validate} y revise \textbf{Validation}. Un modelo válido
 es una entrada coherente; todavía no es una solución.
 \item En \textbf{Run}, seleccione \textbf{QSB Native} y ejecute
 el botón \textbf{Run}; también puede usar \textbf{Solve} en la barra
 de herramientas o \textbf{Solve > Run Current Model} en el menú.
 \item En \textbf{Summary}, compruebe $X_1=18$, $X_2=48$ y objetivo 3780.
 Ambas restricciones se cumplen con igualdad. Revise también el estado de
 resolución y los metadatos antes de interpretar el número como óptimo.
 \item Exporte el modelo y la solución por separado. Vuelva a abrir el modelo
 exportado para conservar una entrada reproducible.
\end{enumerate}
El menú \textbf{Samples} permite explorar ejemplos internos, pero estos pueden
ser distintos de los archivos del manual. Para reproducir un resultado citado,
use el archivo indicado en su capítulo.

\section{Crear, importar y editar}
\textbf{New Model...} abre el selector de familias y crea un modelo inicial.
\textbf{Open Model...} acepta JSON normalizado y formatos heredados reconocidos.
La familia determina los controles que aparecen. Los editores nativos permiten
trabajar con nombres, tablas y campos tipados; otras variantes usan un formulario
estructurado y el JSON como superficie de intercambio.

El panel \textbf{Definition} define el problema; \textbf{Overview} resume su
familia y estado. \textbf{JSON} permite inspeccionar la representación. Cuando
se ofrece \textbf{Apply JSON to ... Editor}, aplique explícitamente el texto
editado para reemplazar el borrador nativo. Escribir JSON no significa que sus
cambios ya estén incorporados a la tabla. Revise los diagnósticos si falla la
aplicación y conserve el texto que necesita reparar.

Algunas tablas ofrecen \textbf{Copy Table} y \textbf{Paste Table} para intercambiar
rangos con hojas de cálculo. Revise encabezados, orden de columnas, separadores,
unidades y nombres tras pegar. Los campos numéricos malformados no se convierten
silenciosamente en ceros. En los editores que señalan dimensiones importadas
incompletas, revise las celdas y los valores mostrados antes de aceptar
\textbf{Use Displayed Dimensions}.

\section{Validar, ejecutar e interpretar estados}
La aplicación distingue modelo vacío, en edición, en validación, inválido y
válido. La ejecución tiene estados independientes: \textbf{Not run},
\textbf{Validating}, \textbf{Solving}, \textbf{Solved} y \textbf{Failed}.
Un estado Solved debe leerse junto a la condición matemática que informa la
familia: un método puede devolver una aproximación, una solución local o una
política heurística en lugar de un óptimo global.

\textbf{Validation} muestra errores y advertencias del modelo. Corrija los
 errores antes de resolver; una advertencia exige juicio, aunque permita seguir.
\textbf{Run Details} permite revisar el backend y la información de ejecución.
Cuando una ejecución larga muestra progreso, espere a que termine antes de
repetir Run. No hay un control universal de cancelación; las rutas asíncronas
de algunos editores evitan que una respuesta antigua sustituya un borrador
nuevo. Esto no equivale a cancelar el cálculo en segundo plano.

\begin{description}
 \item[QSB Native] Métodos educativos deterministas. La exactitud y los límites
 dependen de la familia y se explican en sus capítulos.
 \item[Validate only] Comprueba la entrada y produce diagnósticos; no calcula
 decisiones ni valores óptimos.
 \item[External solver] Integraciones disponibles sólo para determinadas familias.
 La disponibilidad depende de la configuración y del solver instalado.
 No interprete la presencia de la opción como soporte externo universal.
\end{description}
Para LP/ILP y QP, consulte las condiciones de HiGHS en el capítulo de
optimización. Las demás familias pueden indicar que el backend externo no está
disponible. Cambie a QSB Native o Validate only según el trabajo que quiera hacer.

\section{Guardar modelos y resultados}
El botón \textbf{Export} de la barra prepara el JSON a partir del borrador
actual. Si un campo nativo no puede convertirse al modelo, QSB muestra
validación en vez de exportar datos antiguos. El menú
\textbf{Model > Export Model JSON...} abre la exportación del texto de modelo;
para asegurar que los campos recientes estén sincronizados, use primero
Validate o el botón Export de la barra.
\textbf{Export Solution JSON...} guarda el resultado existente y se habilita
cuando hay contenido de solución.

Son documentos distintos: el modelo describe qué problema resolver y la
solución registra qué devolvió una ejecución. Los formatos de salida varían
entre familias; algunos incluyen el modelo y otros no. Conserve siempre ambos,
además del comando, backend, versión del software y supuestos de su estudio.
No use un documento de solución como si fuera una entrada de modelo.

\section{Teclado y alternativas a los gráficos}
Los atajos definidos por la aplicación son:
\begin{center}
\begin{tabular}{ll}
\toprule
Acción & Teclas\\
\midrule
New Model & Cmd+N\\
Open Model & Cmd+O\\
Export Model JSON & Cmd+S\\
Export Solution JSON & Cmd+Shift+S\\
Validate Model & Cmd+Shift+V\\
Run Current Model & Cmd+R\\
Solve ILP & Cmd+I\\
\bottomrule
\end{tabular}
\end{center}
Los botones se habilitan según la familia y el estado. Solve ILP y Solve LP
Relaxation son acciones específicas de programación lineal/entera; la relajación
no impone integralidad. Algunos editores usan Cmd+Shift+N para añadir filas.
Use Tab para desplazarse por controles nativos y los menús para localizar las
acciones; el detalle del foco depende del editor.

Las tablas y el JSON permiten consultar los datos sin depender del color de un
gráfico. Los resultados nativos incorporan etiquetas y valores accesibles en
varias familias, pero esto no acredita por sí solo una revisión completa con
VoiceOver. Para un análisis reproducible, conserve las tablas numéricas. La
ventana macOS tiene mínimo de 980 por 620 puntos; use las áreas desplazables y
cierre el Inspector cuando necesite más anchura.

\section{Terminal y JSON normalizado}
Desde la raíz del repositorio, SwiftPM construye y ejecuta la herramienta:
\begin{lstlisting}
swift run qsb
swift run qsb validate manual/examples/lp_sample.json
swift run qsb solve-json manual/examples/lp_sample.json --backend native
\end{lstlisting}
Sin argumentos, \texttt{qsb} muestra el catálogo de comandos. Tras
\texttt{swift build}, puede ejecutar \texttt{.build/debug/qsb} directamente.
Para guardar la salida y separar errores:
\begin{lstlisting}
.build/debug/qsb solve-json manual/examples/lp_sample.json \
  --backend native > /tmp/qsb-lp-solution.json 2> /tmp/qsb-lp-errors.txt
\end{lstlisting}
El JSON usa punto decimal, cadenas entre comillas dobles y \texttt{null} para
los campos opcionales que lo admiten. No incluye comentarios. Cada familia
posee un esquema propio: LP usa un objeto directo, mientras muchas otras usan
un envoltorio con \texttt{kind} y \texttt{model}. Copie la estructura del ejemplo
de la familia, manteniendo alineados nombres, vectores y matrices.

\texttt{solve-json} detecta la familia de una entrada normalizada.
\texttt{validate} acepta entradas normalizadas o heredadas.
\texttt{solve} está dirigido a modelos heredados; no sustituya indiscriminadamente
un comando por otro. Los capítulos ofrecen los comandos específicos de cada
familia, como \texttt{solve-dp-json} o \texttt{solve-decision-json}.
\texttt{validate-json} es específico de LP/ILP, no un validador universal.

Para importar una entrada heredada propia, la ruta general es:
\begin{lstlisting}
qsb import-legacy-json /ruta/al/modelo > /tmp/qsb-model.json
qsb validate /tmp/qsb-model.json
qsb solve-json /tmp/qsb-model.json --backend native
\end{lstlisting}
Aquí \texttt{qsb} presupone que el ejecutable está en PATH; puede sustituirlo
por \texttt{swift run qsb}. Los archivos originales deben ser accesibles
localmente y usados conforme a sus permisos. Los ejemplos del manual son
independientes de materiales propietarios.

\section{Diagnóstico de problemas}
\begin{description}
 \item[JSON ilegible] Revise sintaxis, nombres de claves y tipos; un error de
 decodificación puede ocurrir antes de generar diagnósticos estructurados.
 \item[Modelo inválido] Consulte \texttt{isValid}, gravedad, código, mensaje y
 ruta del diagnóstico cuando la salida los incluya. Corrija la entrada y vuelva
 a validar.
 \item[Sin solución factible] La estructura puede ser válida y, aun así,
 impedirse satisfacer las restricciones o demandas. Revise límites y balances.
 \item[Backend no disponible] Compruebe qué familias lo soportan y la
 configuración. Validate only no proporciona un resultado alternativo de solver.
 \item[Resultado inesperado] Revise unidades, signo del objetivo, orden de
 estados, restricciones, opciones y exactitud. Compare con cálculos sencillos
 antes de cambiar de algoritmo.
\end{description}
La CLI termina con éxito cuando completa el comando y con fallo ante errores
lanzados, incluido uso incorrecto. Algunas validaciones retornan normalmente un
informe con \texttt{isValid: false}; el código de salida por sí solo no prueba
que el modelo sea válido. Un flujo automático debe revisar también el documento
de validación y el estado del resultado. La ausencia de errores de sintaxis
no demuestra que el modelo represente bien el sistema real.
